\chapter{Rigid-Body Motion}
\label{ch:rigid}

Everything the arm does is described by rotations and translations of rigid bodies. This
chapter is the vocabulary; chapters~\ref{ch:fk} and \ref{ch:jac} are the grammar. Reference:
Modern Robotics ch.~3 (Lynch \& Park); read this chapter alongside it.

\section{Rotations}

A rotation matrix $R \in SO(3)$ is a $3\times3$ matrix satisfying
\[
  R^\top R = I, \qquad \det R = +1 .
\]
Its columns are the axes of the rotated frame expressed in the reference frame — that reading
solves most sign confusions. Rotations compose by multiplication and \emph{do not commute}:
rotating $90^\circ$ about $\hat z$ then $\hat x$ lands differently than $\hat x$ then $\hat z$.
Order is meaning: $R_{ab}R_{bc} = R_{ac}$ (subscript chains cancel like units).

\paragraph{Representations we will actually touch.}
\begin{itemize}
  \item \textbf{Rotation matrix} — what TF2 and all math uses internally; 9 numbers, 3 DoF.
  \item \textbf{Axis--angle} $(\hat\omega, \theta)$ — Rodrigues' formula reconstructs $R$:
  \begin{equation}
    R = I + \sin\theta\,[\hat\omega] + (1-\cos\theta)\,[\hat\omega]^2 ,
    \label{eq:rodrigues}
  \end{equation}
  where $[\hat\omega]$ is the skew-symmetric matrix of $\hat\omega$
  ($[\hat\omega]x = \hat\omega \times x$). This \emph{is} the exponential map:
  $R = e^{[\hat\omega]\theta}$.
  \item \textbf{Unit quaternion} $q = \big(\cos\tfrac\theta2,\; \hat\omega\sin\tfrac\theta2\big)$
  — 4 numbers, no singularities, what ROS messages carry (\texttt{x,y,z,w} order in
  \texttt{geometry\_msgs}, $w$ last — a classic bug source). $q$ and $-q$ are the same rotation
  (double cover).
  \item \textbf{Euler/RPY} — human-readable, gimbal-locked, used only at config-file boundaries
  (URDF \texttt{rpy=}). Never do math in RPY.
\end{itemize}

\begin{example}[Quaternion of a $90^\circ$ yaw]
$\hat\omega = (0,0,1)$, $\theta = \pi/2$:
$q = (\cos 45^\circ,\, 0,0,\sin 45^\circ) = (\tfrac{\sqrt2}{2}, 0, 0, \tfrac{\sqrt2}{2})$.
In ROS message order: $x{=}0,\, y{=}0,\, z{=}0.7071,\, w{=}0.7071$.
\end{example}

\section{Homogeneous transforms}

Pose = rotation + translation, packed as $T \in SE(3)$:
\[
  T = \begin{bmatrix} R & p \\ 0 & 1 \end{bmatrix}, \qquad
  T^{-1} = \begin{bmatrix} R^\top & -R^\top p \\ 0 & 1 \end{bmatrix} .
\]
Composition chains frames: $T_{ab}\,T_{bc} = T_{ac}$. This subscript-cancellation rule is the
whole of TF2: \texttt{lookup\_transform(a, c)} multiplies the chain of stored transforms for
you.

\begin{keyresult}
Never invert by transposing the whole $T$; only $R$ transposes. The translation picks up
$-R^\top p$.
\end{keyresult}

\begin{example}[Wrist camera to base]
The wrist camera sees an ArUco tag at $T_{\text{cam,tag}}$. The arm reports
$T_{\text{base,tool0}}$; hand--eye calibration gave $T_{\text{tool0,cam}}$. Where is the tag
for the planner?
\[
  T_{\text{base,tag}} \;=\; T_{\text{base,tool0}}\; T_{\text{tool0,cam}}\; T_{\text{cam,tag}} .
\]
Every panel-servicing skill begins with exactly this triple product. A $2$\,mm error in the
middle factor (calibration) shifts every tag estimate by $2$\,mm --- which is why
chapter~B1 (hand--eye) exists.
\end{example}

\section{Twists and screws (the 60-second version)}

Angular and linear velocity pack into a twist $\mathcal V = (\omega, v) \in \mathbb R^6$. A
unit screw axis $\mathcal S = (\hat\omega,\, v)$ with $v = -\hat\omega \times q$ (for a
revolute joint through point $q$) generalizes ``rotation about a line.'' Matrix exponentials of
screws produce rigid displacements: $T = e^{[\mathcal S]\theta}$. That single fact powers the
next chapter's forward kinematics. For deeper treatment: MR~ch.~3.3; for the Lie-group view,
Sol\`a et al.\ (arXiv:1812.01537) --- read it \emph{after} this clicks numerically.

\section{Problems}

\begin{problem}
Show that $R_z(90^\circ)R_x(90^\circ) \neq R_x(90^\circ)R_z(90^\circ)$ by tracking where the
unit vector $\hat x$ lands under each. (No matrices needed --- rotate a pen.)
\end{problem}
\begin{answer}
First ordering (apply $R_x$ first, then $R_z$... careful: in a fixed frame, the \emph{left}
matrix acts last): $R_z R_x$ sends $\hat x \to R_x \hat x = \hat x \to R_z \hat x = \hat y$.
The other ordering sends $\hat x \to \hat y \to$ (rotation about $\hat x$) $\to \hat z$.
$\hat y \neq \hat z$: non-commutative.
\end{answer}

\begin{problem}
Using \eqref{eq:rodrigues}, compute $R$ for $\hat\omega = (0,0,1)$, $\theta = 90^\circ$, and
verify it maps $\hat x \to \hat y$.
\end{problem}
\begin{answer}
$[\hat\omega] = \begin{bmatrix}0&-1&0\\1&0&0\\0&0&0\end{bmatrix}$, $\sin\theta = 1$,
$1-\cos\theta = 1$, $[\hat\omega]^2 = \mathrm{diag}(-1,-1,0)$. Sum:
$R = \begin{bmatrix}0&-1&0\\1&0&0\\0&0&1\end{bmatrix}$; first column $= \hat y$. \qedhere
\end{answer}

\begin{problem}
$T_{\text{base,tool0}}$ has $R = R_z(90^\circ)$ and $p = (0.4, 0.1, 0.5)$. Compute
$T^{-1}$ and state what frame relationship it represents.
\end{problem}
\begin{answer}
$R^\top = R_z(-90^\circ)$;
$-R^\top p = -R_z(-90^\circ)(0.4,0.1,0.5) = -(0.1, -0.4, 0.5) = (-0.1, 0.4, -0.5)$.
It is $T_{\text{tool0,base}}$: the base pose as seen from the tool.
\end{answer}

\begin{problem}[transfer]
Your teammate reports the gripper quaternion as $(0.7071, 0, 0, 0.7071)$ ``in ROS order.''
What are the two possible rotations they might mean, and how do you disambiguate?
\end{problem}
\begin{answer}
If $(x,y,z,w)$: $x = 0.7071, w = 0.7071$ → $90^\circ$ about $\hat x$. If they wrote $(w,x,y,z)$:
$90^\circ$ about $\hat z$. Disambiguate by convention-checking against a known pose (or just
\texttt{ros2 topic echo} and compare with RViz). Adopted team rule: always name the convention
when writing quaternions down.
\end{answer}
